\chapter{Mathematical Foundations}
\label{ch:math}

\section{Ontologies as Riemannian Manifolds}

\begin{definition}[Ontology Manifold]
An ontology $\onto$ is a triple $(V, E, \Sigma)$ where:
\begin{itemize}
    \item $V$ is the set of concept nodes
    \item $E \subseteq V \times V$ is the set of semantic relationships
    \item $\Sigma: V \to \mathcal{L}$ maps concepts to logical specifications
\end{itemize}
We endow $\onto$ with a Riemannian metric $g_{ij}$ where $g_{ij} = \expect{\nabla \Sigma_i, \nabla \Sigma_j}$.
\end{definition}

\section{The Generation Fiber Bundle}

Code generation can be formalized as a fiber bundle $\pi: E \to \onto$ where:
\begin{itemize}
    \item Base space: ontology manifold $\onto$
    \item Total space: code manifold bundle $E$
    \item Fiber: $\pi^{-1}(v)$ = all code implementations of concept $v$
    \item Projection: $\pi$ maps code back to originating concept
\end{itemize}

\begin{theorem}[Geodesic Generation]
The optimal code generation path is a geodesic in the total space $E$ with respect to the induced metric.
\end{theorem}

\section{Christoffel Symbols for SPARQL Queries}

SPARQL graph pattern matching computes the Christoffel symbols $\Gamma^k_{ij}$ that define parallel transport in the ontology manifold:

\begin{equation}
\Gamma^k_{ij} = \frac{1}{2} g^{kl} \left( \frac{\partial g_{jl}}{\partial x^i} + \frac{\partial g_{il}}{\partial x^j} - \frac{\partial g_{ij}}{\partial x^l} \right)
\end{equation}

This explains why SPARQL is a natural language for ontology traversal: it computes the connection coefficients!

\section{Information Density Tensor}

\begin{definition}[Information Density Tensor]
The information density at point $p \in \onto$ is the Ricci tensor:
\begin{equation}
\text{Ric}_{ij} = R^k_{ikj} = \frac{\partial \Gamma^k_{ij}}{\partial x^k} - \frac{\partial \Gamma^k_{ik}}{\partial x^j} + \Gamma^k_{ij}\Gamma^l_{kl} - \Gamma^k_{il}\Gamma^l_{kj}
\end{equation}
\end{definition}

High curvature regions (many interconnected concepts) have high generative capacity.
